\section{State of the Art}\label{sec:sota}


\subsection{Neutral Atom Readout}


Conventional readout methods sum pixel intensities within a region surrounding each trap site and compare the result to a threshold. The summation weights may be uniform (square filter) or follow a two-dimensional Gaussian profile centered on the atom, with the width determined by fitting to averaged images~\cite{kent2025}. Winklmann et al.\ compare these filters with Wiener and Richardson--Lucy deconvolution, projection-based state reconstruction~\cite{wei2023}, and both local and global least squares fitting approaches. They also derive a Cramér–Rao bound for achievable precision and find that a global non-linear least squares fit, as introduced by Wei~\cite{wei2023}, offers the highest precision, at the cost of computational resources~\cite{winklmann2024}.


Crosstalk becomes the dominant problem when the atoms are close together. Impertro et al.\ study a cesium quantum gas microscope with a lattice spacing of \SI{383.5}{\nano\metre} and a Rayleigh resolution of \SI{850}{\nano\metre}, i.e.\ a resolution-to-spacing ratio of $\beta = 2.2$. Their convolutional autoencoder is trained without labels directly on experimental images, takes $16\times16$ lattice sites ($256\times256$ pixels) as input, reaches reconstruction fidelities of \SI{96}{\percent}, and reconstructs several thousand sites in less than one second~\cite{impertro2023}.


Supervised neural readout for tweezer arrays is introduced by Phuttitarn et al.\ on a $3\times3$ cesium array. Their CNN site is a five-layer classifier (three convolutional and two fully connected layers) that receives a $10\times10$ pixel crop of one site and has about $3.6\times10^{5}$ parameters. Their CNN array is a six-layer network that takes the entire image as input and has one fully connected head per qubit. It has about $7.5\times10^{7}$ parameters, roughly 210 times as many as the CNN site. Compared with a Gaussian threshold, at \SI{9}{\micro\metre} spacing, where crosstalk is negligible, the error is reduced by up to \SI{56}{\percent} (site) and \SI{32}{\percent} (array). At \SI{5}{\micro\metre} spacing, where crosstalk matters, the reductions are up to \SI{32}{\percent} (site) and \SI{43}{\percent} (array), and the array model lowers the crosstalk correlation by up to \SI{78.5}{\percent}. The authors consider the CNN site small enough for use on a neural network accelerator or FPGA, but consider a real-time implementation of the CNN array challenging~\cite{phuttitarn2023}.


Kent et al.\ address computational cost using the same dataset. Their MF site model is a matched filter in which pixels within a square region of width $s$ around each site are multiplied by learned weights, summed, and thresholded. Their MF array model extends this approach by appending the mean pixel value from the region surrounding each neighboring site to the feature vector, allowing detection of light spillover from adjacent bright sites. Compared to the Gaussian threshold, the error is reduced by up to \SI{32}{\percent} (MF site) and \SI{43}{\percent} (MF array). The MF array requires approximately two orders of magnitude fewer parameters and four orders of magnitude fewer multiplications than the CNN array, at the cost of a \SI{3.5}{\percent} higher error~\cite{kent2025}. The MF site, MF array, and CNN site models serve as baselines for this report.


Two recent works change the input or the classifier's conditions. Mude et al.\ place a fully convolutional denoising network (4.7 million parameters) before a lightweight classifier. This allows reliable classification at up to $1.6\times$ shorter readout times and, compared with state-of-the-art CNN readout, up to $1.77\times$ shorter error correction cycles and up to $35\times$ lower logical error rates for bivariate bicycle codes; the denoiser is evaluated on a GPU~\cite{mude2025}. Su et al.\ keep the CNN site as a base and feed the exposure time as an additional scalar, which is concatenated to the feature vector after the first fully connected layer. This adds only 64 parameters and allows a model to generalize across exposure times, including an exposure time left out of the training set~\cite{su2026}. The model conditioning in \cref{sec:network} follows the same idea, with information about the neighboring sites as the conditioning input.


\subsection{Neural Networks on FPGAs}


Su et al.\ implement their adaptive CNN site model on a ZCU208 board as a fully pipelined design, using $20\times20$ pixel patches and 8-bit fixed-point arithmetic. This implementation achieves classification accuracy within \SI{0.05}{\percent} of the floating-point model. The measured latency is \SI{9.85}{\micro\second} per patch at \SI{500}{\mega\hertz}, compared to \SI{298.85}{\micro\second} on an Intel i9-14900K CPU and \SI{357.12}{\micro\second} on an NVIDIA RTX 4090 GPU. The design uses \SI{6.31}{\percent} of the look-up tables (LUTs), \SI{18.73}{\percent} of the digital signal processing (DSP) blocks, and \SI{45.46}{\percent} of the block RAM (BRAM). The reported latency measures the interval from patch entry into the accelerator to the classification result, excluding image transmission and preprocessing~\cite{su2026}. This work provides the closest reference for the present hardware target.


The accelerator by Winklmann et al.\ uses the entire camera image and implements the projection-based state reconstruction of Wei~\cite{wei2023} with a $31\times31$ point-spread-function kernel in high-level synthesis (HLS) compatible C++ (Vitis HLS 2024.2) on a ZCU216 board at \SI{100}{\mega\hertz}, with parallel vector units, an adder tree, and a prefetching mechanism. A $256\times256$ pixel image of a $10\times10$ array is analyzed in \SI{115}{\micro\second}, which is $34.9\times$ faster than the original CPU implementation and $6.3\times$ faster than their optimized CPU version. For a $40\times40$ array ($1024\times1024$ pixels) the run time grows to \SI{1.825}{\milli\second} and the speedups drop to $20.6\times$ and $2.8\times$. The design uses \SI{25.71}{\percent} of the LUTs, \SI{15.46}{\percent} of the flip-flops, \SI{10.46}{\percent} of the DSPs and \SI{6.20}{\percent} of the BRAM, independent of the array size. The test images are simulated using the simulator in \cite{winklmann2023} and preloaded into double data rate (DDR) memory, rather than captured from a live camera~\cite{winklmann2026accelerator}.

\medskip

In summary, current research reveals a gap in the literature. The neural network implemented on an FPGA processes one site patch at a time, and its latency excludes preprocessing~\cite{su2026}. The models that use the neighboring sites, CNN array and MF array, are evaluated only in software~\cite{phuttitarn2023,kent2025}, and the denoising approach runs on a GPU~\cite{mude2025}. This report integrates most of these approaches into a compact, neighbor-aware CNN that can be deployed in parallel to increase prediction throughput, as detailed in \cref{sec:network,sec:quant,sec:fpga}.
